\subsection{Evaluation experimentale via les API reelles}

Un scenario experimental bout en bout a ete execute via les API reelles
d'AdaptiveEngine afin de verifier la coherence fonctionnelle des recommandations
dans un cas controle. Le scenario cree un cours de test, un graphe de prerequis,
des ressources, des evaluations, des travaux pratiques, puis simule un test de
positionnement et un echec formatif sur le concept \textit{Conditions}.

L'oracle pedagogique attendu est defini a partir des regles du moteur : si le
concept \textit{Conditions} est echoue, il doit etre recommande en remediation ;
le concept \textit{Variables}, deja maitrise, ne doit pas etre priorise comme
lacune principale ; les concepts \textit{Boucles}, \textit{Fonctions} et
\textit{Tableaux} ne doivent pas etre prioritaires avant la remediation de
\textit{Conditions}.

Les recommandations retournees par le moteur sont :
\begin{verbatim}
["ae-real-concept-conditions", "ae-real-concept-boucles", "ae-real-concept-fonctions", "ae-real-concept-tableaux", "ae-real-concept-variables"]
\end{verbatim}

Les recommandations attendues sont :
\begin{verbatim}
["ae-real-concept-conditions"]
\end{verbatim}

\begin{table}[H]
\centering
\caption{Metriques de recommandation calculees sur le scenario API reel}
\begin{tabular}{lccc}
\hline
Scenario & Precision@3 & Recall@3 & nDCG@3 \\
\hline
Conditions echoue & 0.333 & 1.000 & 1.000 \\
\hline
\end{tabular}
\end{table}

Cette evaluation reste limitee : elle valide la coherence du classement dans un
scenario controle execute via les microservices reels, mais ne prouve pas une
amelioration effective de l'apprentissage. Une validation pedagogique robuste
necessiterait plusieurs apprenants, plusieurs cours et un protocole experimental
a plus grande echelle.
